\section{A factor risk model from public data}
\label{sec:model}
Replicating \eqref{eq:qp} requires a $\Sig$. This section builds one in the
conventional multi-factor idiom, from descriptors a reader can compute from
free sources, and then builds three alternatives to it so that
Section~\ref{sec:estimator} can ask whether the choice matters. Everything is
weekly: exposures are measured on Fridays, factor returns span Friday to
Friday, and the last available trading day substitutes when a Friday is a
holiday.

\subsection{Exposures}
Six style descriptors are computed for every name in the weekly cross-section:
the logarithm of market capitalisation; the beta of daily excess returns to the
cap-weighted parent over a trailing 252 days, shrunk toward one; the standard
deviation of the residual from that same regression; the cumulative return over
252 days skipping the most recent 21, which is the usual construction that
avoids one-month reversal; book equity over market capitalisation, with book
equity taken point-in-time as Section~\ref{sec:data} describes; and the
logarithm of average daily dollar volume over a trailing 63 days.

Beta is shrunk following \citet{blume1971assessment},
\begin{equation}
 \hat\beta_i^{\,\rm shrunk} = 0.67\,\hat\beta_i^{\,\rm raw} + 0.33,
 \label{eq:blume}
\end{equation}
with coefficients summing to one so the parent portfolio retains a beta of one.
The two-thirds weighting is a conventional choice rather than a value estimated
on this sample, and nothing in the results turns on the exact split.

Each raw descriptor is winsorised at its 1st and 99th cross-sectional
percentiles and then standardised twice over, to a capitalisation-weighted mean
of zero and an equal-weighted standard deviation of one:
\begin{equation}
 X_{ik} = \frac{d_{ik} - \sum_j b_j d_{jk}}
               {\operatorname{sd}_{\rm ew}\!\left(d_{\cdot k} -
                \textstyle\sum_j b_j d_{jk}\right)}.
 \label{eq:standardise}
\end{equation}
The asymmetry is deliberate and it is not cosmetic. Centring on the
cap-weighted mean makes the parent index's exposure exactly zero on every
style, which is what allows the $\pm0.25$ band in \eqref{eq:qp} to be read as a
deviation from the parent rather than from an arbitrary origin. Scaling by the
equal-weighted standard deviation keeps the unit of that band from being
dictated by the largest few names. A band stated in standard deviations is only
as meaningful as the standardisation underneath it, so a replication that
standardises differently from the provider is applying a different constraint
while appearing to apply the same one.

Standardisation uses only the names in that week's cross-section. Standardising
across the pooled history would let later information set the scale for earlier
weeks, which is look-ahead through a side door.

\subsection{Factor returns and an identification problem}
Weekly factor returns come from a cross-sectional regression of returns on the
exposures known at the start of the week,
\begin{equation}
 r_{i,t} = f^{\rm mkt}_t
         + \sum_j \Xind_{ij,t-1} f^{\rm ind}_{j,t}
         + \sum_k \Xsty_{ik,t-1} f^{\rm sty}_{k,t}
         + u_{i,t},
 \label{eq:csreg}
\end{equation}
fitted by weighted least squares with weights proportional to $\sqrt{b_i}$. The
square root is the Barra convention and sits between two bad alternatives: a
cap-weighted fit is effectively a regression on thirty mega caps, an
equal-weighted fit hands the estimate to the noisiest small names. Exposures
are dated $t-1$ and returns span $t-1$ to $t$. Several descriptors are built
from trailing returns, so fitting contemporaneous exposures would let the
regression explain part of the return with itself.

Equation~\eqref{eq:csreg} is rank deficient as written. Industry membership is
exhaustive and exclusive, so the industry dummies sum to the intercept for
every row and the market factor is not separately identified. The textbook
remedy is to drop one industry column, and it works numerically while quietly
changing what every coefficient means: the intercept stops being the market
return and becomes the omitted industry's return, and each remaining industry
coefficient becomes a difference against it. A market factor defined that way
is not comparable with anything built the conventional way.

We impose instead the economically meaningful side condition that cap-weighted
industry returns sum to zero,
\begin{equation}
 \sum_j c_{j,t}\, f^{\rm ind}_{j,t} = 0,
 \qquad c_{j,t} = \text{industry } j\text{'s share of total capitalisation},
 \label{eq:zerosum}
\end{equation}
by reparametrising rather than by deleting. Writing $f^{\rm ind} = Rg$ with $R$
chosen so that $c^{\top}R = 0$, fitting the unconstrained regression in $g$ and
mapping back leaves the intercept as the cap-weighted market return and each
industry return as a deviation from it. Appendix~\ref{app:identification} gives
$R$ and shows that the two approaches produce different market factors rather
than the same one in different clothing. Industries holding fewer than five
names in a given week are merged into the catch-all for that week, since an
industry return estimated from three stocks is an average of three stocks.

\subsection{Four covariance estimators}
The structural estimator assembles
\begin{equation}
 \Sig_t = X_{t}\,F_{t}\,X_{t}^{\top} + \Delta_{t},
 \label{eq:structural}
\end{equation}
with $F$ the factor covariance and $\Delta$ diagonal. $F$ is estimated by
exponentially weighted moving average with separate half-lives for volatilities
and correlations, 17 weeks and 50 weeks respectively, over a 260-week window.
Volatilities get the shorter memory because they move with regime and have to
react; correlations are more stable and a longer window buys precision at
little cost in staleness. Specific variances use a 26-week half-life and are
shrunk 30\% toward the mean specific variance of the stock's
capitalisation decile, which stops a newly listed or thinly traded name
contributing an estimation artefact the optimiser would happily exploit.

Three alternatives run on the same weekly returns over a common 156-week
window, so the comparison is like with like. The sample covariance is the
baseline, and with roughly five hundred names and 156 observations it is
singular by construction; that is not a defect in the baseline so much as the
baseline's defining problem, and Section~\ref{sec:estimator} reports what it
does anyway. Constant-correlation shrinkage follows
\citet{ledoit2003improved} and \citet{ledoit2004honey}, shrinking toward a
target in which every pairwise correlation equals the cross-sectional average,
with the intensity computed analytically rather than tuned. It is worth
stating which estimator this is: the shrinkage most libraries supply under the
Ledoit--Wolf name targets scaled identity, a different estimator from a
different paper, and the two are routinely conflated. The fourth extracts
principal components of the return covariance and retains $k$ of them, with $k$
chosen by the \citet{bai2002determining} ICp2 criterion over a search bounded
at twenty rather than fixed at a round number chosen after seeing results.

A long-only minimum-variance portfolio built on US equities from a return-based
covariance estimate is the design \citet{clarke2006minimum} study directly, and
the objective itself descends from \citet{markowitz1952portfolio}, whom the
index provider's own methodology cites. The four estimators span the trade-off
\citet{briner2008howmuch} describe between
estimation variance and specification error: the sample matrix imposes no
structure and pays in variance, the structural model imposes a great deal and
pays in specification error, and shrinkage and principal components sit between
them by different routes. Whether that trade-off survives contact with a
binding constraint set is the question Section~\ref{sec:estimator} puts to
them.
